\documentclass[11pt,a4paper]{article}
\usepackage[margin=1.5cm]{geometry}
\usepackage{amsmath,amssymb}
\usepackage{circuitikz}
\usepackage{tikz}
\usetikzlibrary{arrows.meta,calc,positioning}

\begin{document}

\section*{Audit Test: 10 Fully Perfected Circuit Schematics}

\subsection*{1. Diagram 1.3: 2082 Baishakh Q1 (Mesh Analysis with $V_x/10$)}
\begin{center}
\begin{circuitikz}[scale=1.0, transform shape, american]
    % Left current source branch: downward current of -10A
    \draw (0,4) to[isource, l_=$-10\,\mathrm{A}$] (0,0);
    
    % Top wire connecting left, middle, and right branches
    \draw (0,4) -- (3.5,4) -- (7,4);
    
    % Middle vertical branch: 1 ohm on top, dependent current source Vx/10 (downward) on bottom
    \draw (3.5,4) to[R, l_=$1\,\Omega$] (3.5,2.0)
          to[cI, l_=$V_x/10$] (3.5,0);
          
    % Bridge horizontal resistor with Vx polarity (+ on left, - on right)
    \draw (3.5,2.0) to[R, l=$3\,\Omega$, v=$V_x$] (7,2.0);
    
    % Right vertical branch: 2 ohm on top, 5 ohm on bottom
    \draw (7,4) to[R, l=$2\,\Omega$] (7,2.0)
          to[R, l=$5\,\Omega$] (7,0);
          
    % Bottom wire
    \draw (0,0) -- (3.5,0) -- (7,0);
\end{circuitikz}
\end{center}

\vspace{0.8cm}
\subsection*{2. Diagram 1.4: 2081 Ashwin Q1 (Nodal Analysis with $5\mathrm{A}, 3\Omega, 4V_x, 1\Omega$)}
\begin{center}
\begin{circuitikz}[scale=1.0, transform shape, american]
    % Ground reference line
    \draw (0,0) -- (8.5,0);
    \draw (4.25,0) node[ground]{};
    
    % Left side: 5A downward source and 3-ohm resistor as distinct parallel branches to node V1
    \draw (0,2.5) to[isource, l_=$5\,\mathrm{A}$] (0,0);
    \draw (0,2.5) -- (1.8,2.5);
    \draw (1.8,0) to[R, l=$3\,\Omega$] (1.8,2.5);
    \draw (1.8,2.5) node[circ]{} node[above left]{$V_1$};
    
    % Top bridge between V1 and V2: 2-ohm resistor
    \draw (1.8,2.5) -- (1.8,4.0) to[R, l=$2\,\Omega$] (6.5,4.0) -- (6.5,2.5);
    
    % Middle branch between V1 and V2: 5V source (- on left, + on right) in series with 4-ohm resistor (+ Vx -)
    \draw (1.8,2.5) to[battery, l=$5\,\mathrm{V}$] (4.0,2.5)
          to[R, l=$4\,\Omega$, v=$V_x$] (6.5,2.5);
    \draw (6.5,2.5) node[circ]{} node[above right]{$V_2$};
    
    % Right side: 4Vx downward dependent current source and 1-ohm resistor as distinct parallel branches to node V2
    \draw (6.5,2.5) to[cI, l_=$4V_x$] (6.5,0);
    \draw (6.5,2.5) -- (8.5,2.5) to[R, l=$1\,\Omega$] (8.5,0);
\end{circuitikz}
\end{center}

\newpage
\subsection*{3. Diagram 2.2: 2082 Bhadra Q2 (Inductor Switching with $I, K, R, L$)}
\begin{center}
\begin{circuitikz}[scale=1.05, transform shape, american]
    % Ground line
    \draw (0,0) -- (7.2,0);
    
    % Left independent current source I (pointing up)
    \draw (0,0) to[isource, l=$I$] (0,3.2);
    
    % Top wire
    \draw (0,3.2) -- (2.2,3.2);
    \draw (2.2,3.2) to[short, i=$i$] (4.7,3.2) -- (7.2,3.2);
    
    % Switch K branch (opens at t=0)
    \draw (2.2,0) -- (2.2,0.9) node[circ]{};
    \draw (2.2,2.3) node[circ]{} -- (2.2,3.2);
    \draw[thick] (2.2,0.9) -- (1.7,2.2);
    \draw[-{Latex[length=2.5mm]}] (1.5,1.7) arc (130:75:0.8);
    \node at (2.8, 1.6) {$K\ (t=0)$};
    
    % Resistor R branch
    \draw (4.7,3.2) to[R, l_=$R$] (4.7,0);
    
    % Inductor L branch with i_L and v
    \draw (7.2,3.2) to[L, l=$L$, i=$i_L$, v=$v$] (7.2,0);
\end{circuitikz}
\end{center}

\vspace{0.8cm}
\subsection*{4. Diagram 2.3: 2082 Baishakh Q2 (SPDT Transients with $12\mathrm{V}, 2\Omega, 6\Omega, 2\mathrm{F}, 4\Omega, 3\mathrm{H}$)}
\begin{center}
\begin{circuitikz}[scale=1.05, transform shape, american]
    % Bottom line
    \draw (0,0) -- (7.5,0);
    
    % 12V battery
    \draw (0,0) to[battery2, l=$12\,\mathrm{V}$] (0,3.0) -- (1.2,3.0);
    
    % SPDT Switch terminals
    \draw (1.2,3.0) node[circ]{} node[above]{$a$};
    \draw (1.2,1.5) node[circ]{} node[left]{$b$};
    \draw (2.4,2.25) node[circ]{};
    
    % Switch blade throwing from a to b at t=0
    \draw[thick] (2.4,2.25) -- (1.35,2.9);
    \draw[-{Latex[length=2.5mm]}] (1.0,2.6) to[bend right=30] (1.0,1.8);
    \node at (1.8, 3.2) {$t=0$};
    
    % Branch b with 6-ohm resistor
    \draw (1.2,1.5) to[R, l_=$6\,\Omega$] (1.2,0);
    
    % Series 2-ohm resistor from switch blade
    \draw (2.4,2.25) to[R, l=$2\,\Omega$] (4.5,2.25);
    
    % Central shunt: 2F capacitor
    \draw (4.5,2.25) node[circ]{} to[C, l=$2\,\mathrm{F}$] (4.5,0);
    
    % Right branch: series 4-ohm resistor and 3H inductor
    \draw (4.5,2.25) -- (6.8,2.25) -- (6.8, 2.05)
          to[R, l=$4\,\Omega$] (6.8,1.05)
          to[L, l=$3\,\mathrm{H}$] (6.8,0);
\end{circuitikz}
\end{center}

\vspace{0.8cm}
\subsection*{5. Diagram 3.4: 2081 Ashwin Q3 (SPDT Transients with $1\mathrm{V}, 1\Omega, 1\mathrm{H}, 0.5\Omega, 2\mathrm{H}$)}
\begin{center}
\begin{circuitikz}[scale=1.05, transform shape, american]
    % Bottom line
    \draw (0,0) -- (7.5,0);
    
    % 1V DC source and 1-ohm resistor leading to terminal 'a'
    \draw (0,0) to[battery2, l=$1\,\mathrm{V}$] (0,3.0) to[R, l=$1\,\Omega$] (2.2,3.0);
    
    % SPDT Switch terminals
    \draw (2.2,3.0) node[circ]{} node[above]{$a$};
    \draw (3.6,3.0) node[circ]{} node[above]{$b$};
    \draw (2.9,1.8) node[circ]{};
    
    % Switch blade moving from a to b at t=0
    \draw[thick] (2.9,1.8) -- (2.3,2.9);
    \draw[-{Latex[length=2.5mm]}] (2.4,2.6) to[bend left=30] (3.4,2.6);
    \node at (2.9, 3.4) {$t=0$};
    
    % 1H Inductor below the switch blade
    \draw (2.9,1.8) to[L, l_=$1\,\mathrm{H}$, i=$i_L(t)$] (2.9,0);
    
    % Terminal b leading to parallel 0.5-ohm resistor and 2H inductor
    \draw (3.6,3.0) -- (5.2,3.0) to[R, l_=$\frac{1}{2}\,\Omega$] (5.2,0);
    \draw (5.2,3.0) -- (7.0,3.0) to[L, l=$2\,\mathrm{H}$, v=$V(t)$] (7.0,0);
\end{circuitikz}
\end{center}

\newpage
\subsection*{6. Diagram 4.1: 2083 Baishakh Q4 (Laplace Parallel RLC Transients)}
\begin{center}
\begin{circuitikz}[scale=1.05, transform shape, american]
    % Bottom line
    \draw (0,0) -- (6.8,0);
    
    % 4V DC battery
    \draw (0,0) to[battery2, l=$4\,\mathrm{V}$] (0,2.8);
    
    % Switch opening at t=0
    \draw (0,2.8) -- (0.9,2.8) node[circ]{};
    \draw (2.1,2.8) node[circ]{} -- (3.2,2.8);
    \draw[thick] (0.9,2.8) -- (1.9,3.4);
    \draw[-{Latex[length=2.5mm]}] (1.4,2.85) to[bend right=25] (1.4,3.35);
    \node at (1.5, 3.7) {$t=0$};
    
    % 1F capacitor shunt branch
    \draw (3.2,2.8) node[circ]{} to[C, l_=$1\,\mathrm{F}$] (3.2,0);
    
    % Top rail inductor 0.5H and right vertical 1-ohm resistor
    \draw (3.2,2.8) to[L, l=$0.5\,\mathrm{H}$, v=$v_L(t)$] (6.0,2.8)
          to[R, l=$1\,\Omega$, i=$i(t)$] (6.0,0);
\end{circuitikz}
\end{center}

\vspace{0.8cm}
\subsection*{7. Diagram 4.2: 2082 Bhadra Q4 (Laplace Transients with $10\mathrm{V}, 5\Omega, 10\mu\mathrm{F}, 0.1\mathrm{H}$)}
\begin{center}
\begin{circuitikz}[scale=1.05, transform shape, american]
    % Bottom line
    \draw (0,0) -- (6.5,0);
    
    % 10V DC battery and 5-ohm resistor to terminal 'a'
    \draw (0,0) to[battery2, l=$10\,\mathrm{V}$] (0,3.0) to[R, l=$5\,\Omega$] (2.0,3.0);
    
    % SPDT Switch terminals
    \draw (2.0,3.0) node[circ]{} node[above]{$a$};
    \draw (2.0,1.5) node[circ]{} node[left]{$b$};
    \draw (3.2,2.25) node[circ]{};
    
    % Moving blade from a to b at t=0
    \draw[thick] (3.2,2.25) -- (2.15,2.9);
    \draw[-{Latex[length=2.5mm]}] (1.8,2.6) to[bend right=30] (1.8,1.8);
    \node at (2.6, 3.3) {$t=0$};
    
    % Capacitor 10uF at terminal b
    \draw (2.0,1.5) to[C, l_=$10\,\mu\mathrm{F}$] (2.0,0);
    
    % Right branch: 0.1H inductor with i(t)
    \draw (3.2,2.25) -- (5.8,2.25) to[L, l=$0.1\,\mathrm{H}$, i=$i(t)$] (5.8,0);
\end{circuitikz}
\end{center}

\vspace{0.8cm}
\subsection*{8. Diagram 4.4: 2081 Ashwin Q4 (Laplace Transients with $6\mathrm{V}, 2\Omega, 1\Omega, 0.5\mathrm{H}, 1\mathrm{F}$)}
\begin{center}
\begin{circuitikz}[scale=1.05, transform shape, american]
    % Bottom line
    \draw (0,0) -- (6.8,0);
    
    % 6V DC battery and 2-ohm resistor to terminal 'a'
    \draw (0,0) to[battery2, l=$6\,\mathrm{V}$] (0,3.4) to[R, l=$2\,\Omega$] (2.0,3.4);
    
    % SPDT Switch terminals
    \draw (2.0,3.4) node[circ]{} node[above]{$a$};
    \draw (2.0,2.3) node[circ]{} node[left]{$b$};
    \draw (3.2,2.85) node[circ]{};
    
    % Blade moving from a to b at t=0
    \draw[thick] (3.2,2.85) -- (2.15,3.3);
    \draw[-{Latex[length=2.5mm]}] (1.8,3.1) to[bend right=30] (1.8,2.5);
    \node at (2.6, 3.7) {$t=0$};
    
    % Branch b with series 1-ohm resistor and 0.5H inductor (with ample height)
    \draw (2.0,2.3) to[R, l_=$1\,\Omega$] (2.0,1.15) to[L, l_=$0.5\,\mathrm{H}$] (2.0,0);
    
    % Right branch: 1F capacitor
    \draw (3.2,2.85) -- (6.0,2.85) to[C, l=$1\,\mathrm{F}$, v=$v_C(t)$] (6.0,0);
\end{circuitikz}
\end{center}

\newpage
\subsection*{9. Diagram 6.1: 2083 Baishakh Q6 (Two-Port Z-Parameters with $8\Omega, 10\Omega, 20\Omega, 20\Omega$)}
\begin{center}
\begin{circuitikz}[scale=1.05, transform shape, american]
    % Port 1 terminals
    \draw (0,2.8) node[ocirc]{} node[above left]{$+$} to[short, i=$I_1$] (1.0,2.8)
          to[R, l=$8\,\Omega$] (2.8,2.8)
          to[R, l=$10\,\Omega$] (5.4,2.8) -- (7.0,2.8)
          node[ocirc]{} node[above right]{$+$};
    
    % Port 1 and Port 2 voltage annotations
    \node at (-0.6, 1.4) {\large $V_1$};
    \node at (7.6, 1.4) {\large $V_2$};
    
    % Port 2 incoming current I2 (entering port, neatly above)
    \draw[-{Latex[length=3mm]}] (6.8, 3.5) -- node[above]{$I_2$} (5.6, 3.5);
    
    % Shunt branches of 20 ohms
    \draw (2.8,2.8) node[circ]{} to[R, l=$20\,\Omega$] (2.8,0);
    \draw (5.4,2.8) node[circ]{} to[R, l=$20\,\Omega$] (5.4,0);
    
    % Bottom continuous ground rail with port negative terminals
    \draw (0,0) node[ocirc]{} node[below left]{$-$} -- (7.0,0) node[ocirc]{} node[below right]{$-$};
\end{circuitikz}
\end{center}

\vspace{1.2cm}
\subsection*{10. Diagram 6.2: 2082 Bhadra Q5 (Two-Port Z-Parameters with $3\Omega, \pm 2I_a, -j2\Omega, 1\Omega, j5\Omega$)}
\begin{center}
\begin{circuitikz}[scale=1.05, transform shape, american]
    % Port 1
    \draw (0,2.8) node[ocirc]{} node[above left]{$+$} to[short, i=$I_1$] (0.8,2.8)
          to[R, l=$3\,\Omega$] (3.0,2.8);
    \node at (-0.6, 1.4) {\large $V_1$};
    \node at (7.8, 1.4) {\large $V_2$};
    
    % Middle shunt branch: Controlled source \pm 2I_a in series with capacitor -j2\Omega
    \draw (3.0,2.8) node[circ]{} to[cV, l=$\pm 2I_a$] (3.0,1.4)
          to[C, l=$-j2\,\Omega$] (3.0,0);
          
    % Port 2: Series 1-ohm and j5-ohm
    \draw (3.0,2.8) to[R, l=$1\,\Omega$] (4.8,2.8)
          to[L, l=$j5\,\Omega$] (7.2,2.8) node[ocirc]{} node[above right]{$+$};
          
    % Current Ia entering Port 2 (drawn neatly above line)
    \draw[-{Latex[length=3mm]}] (7.0, 3.5) -- node[above]{$I_a$} (5.8, 3.5);
    
    % Bottom continuous reference line
    \draw (0,0) node[ocirc]{} node[below left]{$-$} -- (7.2,0) node[ocirc]{} node[below right]{$-$};
\end{circuitikz}
\end{center}

\end{document}
